\begin{corollary}[Bernstein--Szeg\H{o} inequality applied to $T'/m$]
Let $m\ge1$ be an integer and let $T:\mathbb{R}\to\mathbb{R}$ be a real
trigonometric polynomial of degree at most $m$ with
$\lvert T(t)\rvert\le1$ for every $t\in\mathbb{R}$. Then
\[
  \frac{\bigl(T'(t)\bigr)^2}{m^2}+\frac{\bigl(T''(t)\bigr)^2}{m^4}\le1
  \qquad\text{for every }t\in\mathbb{R},
\]
where $T'$ and $T''$ denote the first and second pointwise derivatives of $T$.
\end{corollary}
\begin{proof}
Since $m\ge1$, the real number $m$ satisfies $m>0$, so $m\ne0$ and all divisions
by powers of $m$ below are legitimate. By \noderef{TrigPolySmooth}, $T$ is
real-analytic on $\mathbb{R}$, in particular infinitely differentiable; so $T'$
exists and is differentiable, and $T''=(T')'$ exists, at every point of
$\mathbb{R}$.

Define
\[
  U(t)=\frac{T'(t)}{m}\qquad(t\in\mathbb{R}).
\]

\emph{Step 1: $U$ is a real trigonometric polynomial of degree at most $m$.} By
\noderef{TrigPolyDeriv} applied to $T$, the pointwise derivative $T'$ is a real
trigonometric polynomial of degree at most $m$. Now apply
\noderef{TrigPolyLinComb} with both functions taken to be $T'$ (each of degree at
most $m$, as just shown) and with the scalars $c=1/m$ and $d=0$: the function
\[
  t\mapsto \frac{1}{m}\,T'(t)+0\cdot T'(t)
\]
is a real trigonometric polynomial of degree at most $m$. For every
$t\in\mathbb{R}$ that function takes the value
$\frac{1}{m}T'(t)+0=\frac{T'(t)}{m}=U(t)$, so it is $U$, and therefore $U$ is a
real trigonometric polynomial of degree at most $m$.

\emph{Step 2: $\lvert U(t)\rvert\le1$ for every $t$.} The hypotheses of
\noderef{BernsteinSzego} hold for the integer $m$ and the function $T$: indeed
$m\ge1$, $T$ is a real trigonometric polynomial of degree at most $m$, and
$\lvert T(t)\rvert\le1$ for every $t$. That theorem therefore gives
\begin{equation}\label{eq:bss-first}
  T(t)^2+\frac{\bigl(T'(t)\bigr)^2}{m^2}\le1
  \qquad\text{for every }t\in\mathbb{R}.
\end{equation}
Since $T(t)^2\ge0$, \eqref{eq:bss-first} yields
$\bigl(T'(t)\bigr)^2/m^2\le1$. As
$U(t)^2=\bigl(T'(t)/m\bigr)^2=\bigl(T'(t)\bigr)^2/m^2$, this says
$U(t)^2\le1$, i.e. $\lvert U(t)\rvert\le1$.

\emph{Step 3: $U'(t)=T''(t)/m$ for every $t$.} We have $U=\frac1m\,T'$ as
functions on $\mathbb{R}$; $T'$ is differentiable at every point (see the
opening paragraph), and the derivative is linear, so $U$ is differentiable with
\[
  U'(t)=\frac1m\,\bigl(T'\bigr)'(t)=\frac{T''(t)}{m}
  \qquad(t\in\mathbb{R}).
\]

\emph{Step 4: applying the inequality to $U$.} By Step 1, $U$ is a real
trigonometric polynomial of degree at most $m$; by Step 2,
$\lvert U(t)\rvert\le1$ for every $t$; and $m\ge1$. So the hypotheses of
\noderef{BernsteinSzego} hold for the integer $m$ and the function $U$, and that
theorem gives
\[
  U(t)^2+\frac{\bigl(U'(t)\bigr)^2}{m^2}\le1
  \qquad\text{for every }t\in\mathbb{R}.
\]
Substituting $U(t)=T'(t)/m$ and, by Step 3, $U'(t)=T''(t)/m$, and using
\[
  U(t)^2=\frac{\bigl(T'(t)\bigr)^2}{m^2},
  \qquad
  \frac{\bigl(U'(t)\bigr)^2}{m^2}
   =\frac{\bigl(T''(t)/m\bigr)^2}{m^2}
   =\frac{\bigl(T''(t)\bigr)^2}{m^4},
\]
this becomes
\[
  \frac{\bigl(T'(t)\bigr)^2}{m^2}+\frac{\bigl(T''(t)\bigr)^2}{m^4}\le1
  \qquad\text{for every }t\in\mathbb{R},
\]
which is the assertion.
\end{proof}
